\section{The primitive row interface}
\label{sec:interface}

Fix \(h\in\Z\setminus\{0\}\) and put
\begin{equation}
 H=\rad|h|.
 \label{eq:H}
\end{equation}
Let \(0<\delta<1/2\) and
\begin{equation}
 R=\lfloor X^{1/2-\delta}\rfloor,
 \qquad V=\lfloor R^{1/2}\rfloor.
 \label{eq:RV}
\end{equation}
The finite Ramanujan model and its residual are
\begin{equation}
 \Lambda_R(n)=\sum_{q\le R}\frac{\mu(q)}{\varphi(q)}c_q(n),
 \qquad r_R(n)=\Lambda(n)-\Lambda_R(n).
 \label{eq:model-residual}
\end{equation}
All target arguments below are positive integers.  We use the standard
convention \(c_q(n)=\sum_{a\bmod q}^{*}\ee{an/q}\).

Let \(L,D,J,Q\ge1\) be dyadic parameters satisfying
\begin{equation}
 Q=LD,\qquad QJ\asymp X,\qquad
 L>\max(2R,2|h|),
 \label{eq:row-geometry-a}
\end{equation}
and
\begin{equation}
 \ell\in[L/2,2L]\ \text{is prime},\qquad
 d\in[D,2D]\cap[1,V].
 \label{eq:row-geometry}
\end{equation}
The integer \(d\) is squarefree, and in the primitive sector
\begin{equation}
 (d,H)=1,\qquad (j,H)=1.
 \label{eq:primitive-mask}
\end{equation}
Put
\begin{equation}
 m=\ell d\asymp Q,\qquad n_m(j)=mj+h.
 \label{eq:m-row}
\end{equation}
Then \((m,H)=1\), and \eqref{eq:primitive-mask} implies
\begin{equation}
 (n_m(j),mH)=1.
 \label{eq:target-coprime}
\end{equation}

TPC-18 reaches row pairs \(\alpha_i=(\ell_i,d_i)\) in the generic set
\begin{equation}
 \ell_1\ne\ell_2,\qquad
 |m_1-m_2|>QX^{-\kappa},\qquad
 (d_1,d_2)\le X^\kappa,
 \label{eq:generic-rows}
\end{equation}
for fixed \(\kappa>0\).  The exact identities below hold before this
pruning, so any row-pair restriction independent of \(j\), including
\eqref{eq:generic-rows}, may be imposed outside them.

For a uniformly smooth compactly supported family
\(F_{\alpha_1,\alpha_2}\), supported where \(j\asymp J\) and both
targets are positive, define the primitive residual correlation
\begin{align}
 \cC_D(\mathscr A)
 :=\sum_{(\alpha_1,\alpha_2)\in\mathscr A}
 z_{\alpha_1,\alpha_2}
 \sum_{\substack{j\ge1\\(j,H)=1}}
 r_R(n_{m_1}(j))r_R(n_{m_2}(j))
 F_{\alpha_1,\alpha_2}(j/J),
 \label{eq:abstract-correlation}
\end{align}
where \(|z_{\alpha_1,\alpha_2}|\ll X^\eps\) absorbs the source weights,
the M\"obius row signs, and bounded dyadic cutoffs.  Uniform smoothness
means that every derivative of every \(F_{\alpha_1,\alpha_2}\) is
bounded independently of \(X\) and the row pair.  This hypothesis is
essential in the Poisson-completion statements below.

The weighted number of individual rows is
\begin{equation}
 \sum_{\ell\asymp L}\Lambda(\ell)
 \sum_{d\asymp D}\mu(d)^2\ll LD\,X^\eps=QX^\eps.
 \label{eq:row-mass}
\end{equation}
Consequently the row-pair--orbit scale is
\begin{equation}
 Q^2J\asymp XQ.
 \label{eq:natural-scale}
\end{equation}
This is the natural size of the TPC-18 obstruction.  Every claim of
negligibility below is measured against \(XQ(\log X)^{-A}\), with a
fixed power margin whenever one is asserted.

\begin{remark}[Smooth separation]
The original TPC-18 weight contains smooth functions of \(mj/X\) and
\(dj/K\), not merely a common function of \(j/J\).  Fourier inversion
in logarithmic coordinates writes such a compactly supported smooth
weight as an absolutely convergent superposition of separated row and
\(j\) weights.  Truncating the dual variables at \(X^\eps\) has
arbitrary power accuracy.  Thus the uniformly smooth family in
\eqref{eq:abstract-correlation} is the stable interface for the exact
identities; no sharp \(j\)-cutoff is being silently completed.
\end{remark}
